\begin{abstract}
We introduce \textbf{OenoBench}, a benchmark of \PLACEHOLDER{final question count, target 5{,}000} multiple-choice and short-answer questions designed to evaluate factual knowledge of large language models (LLMs) across the wine domain --- a discipline that draws together botany, geography, chemistry, regulation, and sensory analysis, and which is uniquely well-suited to benchmark construction because it is supported by formal professional credentials (WSET, Master Sommelier, Master of Wine) whose syllabi externally validate what counts as expert knowledge. OenoBench is built from \textbf{38{,}104 atomic facts} extracted by 35 provenance-verified scrapers across 15 source types, organised into six domain pillars: wine regions, grape varieties, viticulture, winemaking, producers, and wine business. Questions are produced by five complementary generation strategies (fact-to-question, template, comparative, scenario, distractor mining) and distributed across five generator models (Claude, GPT, Gemini, Llama, deterministic templates) to mitigate single-model bias. Every question passes through a nine-agent automated audit organised into four teams (lexical hygiene, tri-judge correctness, deterministic checks, population statistics) and is calibrated against a human-reviewed gold sheet using Cohen's $\kappa$. We report bias diagnostics including a \emph{Self-Preference Score} that quantifies whether models score disproportionately well on questions generated by themselves, and we hold out a 300-question expert-authored control set for ground-truth comparison. \PLACEHOLDER{Headline evaluation results across frontier LLMs: overall accuracy, per-domain accuracy, SPS values, control-set comparison.}
\end{abstract}
